\chapter{可实施特征与实施例}
\label{ch:patent}

本章要判定的量是：前面的判据能否收成同一腔体上三组互不替代的步骤。三组分别回答“参数看见了没有”“效应走哪条路径”“执行器是否被硬测量拦住”。本章不讨论文书格式。

\section{三组步骤的分界}

\paragraph{要判定的量}
一个步骤属于信息、解释还是安全。一个步骤若同时声称三件事，就不可实施：日志无法判断它失败在哪一组。

\paragraph{继承的文献层与不继承的层}
分界继承第 \ref{ch:arch} 章的标记：动作在进入唯一执行器之前带有信息、品质或降级保温的标记。不继承把三组并成一个“智能程序”分数。

信息组的失败是方差或分位没有下降却改变了停止时刻。解释组的失败是路径符号在后验抽样中不稳定却生成了陈述句。安全组的失败是品质估计进入了断电或投影的不等式。三组可以同时运行，但验收分开。

\paragraph{本节不能推出的结论}
分界本身不是步骤已经在器具上执行的证据。

\section{信息组}

\paragraph{要判定的量}
内部热阻是否在离开平衡的两段边界之后被更新；更新无效时停止时刻是否保持。

步骤如下。

\begin{enumerate}
\item 读取表面温度、腔温、湿度和质量。缺失的通道只扩大不确定度，不用穿刺补齐。本次核只读取表面温度和平均含水。
\item 若处于无相变的平衡，不更新内部热阻。命题 \ref{prop:steady} 给出这一步的理由。
\item 信息赤字仍在时，执行先冷凝方向、后高风的两段组合。规格中的时长不超过 20 s。本次核用 150 s 至 190 s 的固定窗代替库上的选择。
\item 仅当后验方差下降时接受新的热阻。否则锁定该次探针，停止时刻保持更新前的值。核里的接受条件还要求新值更接近真值参照。这一条比规格更弱，实施时必须去掉参照，只保留观测。
\item 厚件上方差仍然高于阈值时，不允许用点估计停机。核里的阈值是方差 0.45 且厚度大于 40 mm。
\end{enumerate}

\paragraph{本节不能推出的结论}
固定窗通过了种子 7 的方差下降，不能推出库中其余探针已经被打分。

\section{解释组}

\paragraph{要判定的量}
同一次蒸汽干预能否输出路径 A、B、C 的区间，并且只在符号稳定时输出语句。

步骤如下。

\begin{enumerate}
\item 用当前后验的参数积分结构方程，父节点保持第 \ref{ch:causal} 章的表，不从本次序列里增边。
\item 将蒸汽轨迹替换为候选值，其余执行器保持，重新积分。这是干预，不是条件相关。
\item 对进度和表面含水分别累加，得到 A、B 以及经过结壳的 C。
\item 在后验抽样上重复。符号改变则不生成陈述句，只保留预测区间。
\end{enumerate}

本次核没有这四步。实施顺序上，解释组排在信息组能提供参数之后，但不能反过来改写信息组的停止规则。解释文本不是执行器。

\paragraph{本节不能推出的结论}
没有路径积分时，界面上的三根条不能当作解释组的输出。

\section{安全组}

\paragraph{要判定的量}
投影不等式的自变量是否与品质后验不相交。

步骤如下。

\begin{enumerate}
\item 品质层只产生参考动作。
\item 投影读取漏电、水位或空烧温升、加热器温度、腔压和门。不读取进度估计。
\item 参考动作被投影到安全集内，再写入执行器。
\item 硬件切断与投影串联，不经过产生参考动作的处理器。
\item 置信度不足时，参考动作改为存储的低功率保持，然后仍然经过同一次投影。
\end{enumerate}

本次只落实了加热器温度上的削减公式，而且平衡点低于阈值，削减不动作。断言用另一阈值证明“不随后验改变”。实施时必须把两处阈值收成一套，并补上尚未采样的四路测量。

\paragraph{本节不能推出的结论}
安全组不产生食品安全限值。0.85 仍是设计线。

\section{在同一腔体上的实施例}

\paragraph{要判定的量}
三组步骤是否使用同一套执行器和两套分开的测量，而不需要为信息动作增加机构。

腔体侧已有的执行器是背部或顶部加热、蒸汽入口、热风和排汽。信息组的两段组合只改变它们的时序和目标函数。品质组在探针结束之后使用同一加热器。安全组在写入前改写指令。

食品侧测量是底部质量和门侧热像。穿刺探针若安装，只放在对照通道，不进入信息组的主观测。安全侧是剩余电流、水位、管温、门和腔压，与食品侧分开布线。演示页面上这三组要能同时看见：边界、状态、路径区间、请求与投影。当前页面没有做到，差距在第 \ref{ch:demo} 章。

一条已经跑过的退化实施例是：十六条控制体、表面温度观测、固定时间窗、饱和线性加热、加热器 ODE。它覆盖信息组的时间表和安全组的一个状态变量，不覆盖解释组。

\paragraph{本节不能推出的结论}
退化实施例的中位误差不能写成三组步骤同时工作时的误差。

\section{结果尚未产生时的文档边界}

下列对象保持为规格，不写成结果：多孔介质真值场、集合滤波、Fisher 矩阵的特征值、探针库打分、三条路径的数值、神经算子、预测控制求解器、二次投影在工作点上的动作、S2 的遮挡与缺通道、T2 与 T3。文档可以引用它们的方程和判定标准。文档不可以引用它们的误差、次数或曲线。

已经产生的结果只有：四条命题断言的通过记录，S1 摘要的 19440 行，五张由图文件构成的图，以及种子 7 的一条时间序列。引用时带上核的限制。

\section{实施顺序}

顺序由依赖决定，不由页面需要决定。

\begin{enumerate}
\item 先固定符号和两节点命题，并保留无界面断言。这一步已经有四条通过的记录。
\item 再把热阻更新改成只依赖观测，去掉真值参照。这一步尚未做。
\item 然后才把十六条控制体换成带 GAB 场、各向异性和收缩的真值，并让估计器与真值分离。
\item 真值隔离通过之后，才运行 S2 和 T2。
\item 算子训练使用单独的教师集，不使用 S1 的同一批行做测试。
\item 演示只读取已经通过断言的状态流。着色器里不写第二套差分。
\end{enumerate}

\paragraph{本节不能推出的结论}
顺序靠后的步骤尚未执行，不能用靠前步骤的数字提前填上。
